\section{Parity changes the critical exponent of Ehrhart zeros}
\label{sec:ehrhart-critical}

\subsection{The family and what is being resolved}
For $d\ge2$, let $G_{d,M}$ be the complete graph on $d+1$ vertices with a matching of
$M$ edges deleted. For integer $0\le M\le\lfloor(d+1)/2\rfloor$, let
$E_{d,M}(s)$ be the Ehrhart polynomial of its symmetric edge polytope.
The lattice is the integer zero-sum hyperplane, and the polytope is the
convex hull of $\pm(e_i-e_j)$ over the remaining edges. Write $H_d=E_{d,0}$
and let $C_d$ denote the Ehrhart polynomial of the $d$-dimensional
cross-polytope. The exact identity from the incoming multipartite
report~\cite{EhrhartIncoming} is
\begin{equation}\label{ae:exact}
 E_{d,M}(s)=H_d(s)-M D_d(s),\qquad
 D_d(s)=\frac{C_d(s)-C_{d-2}(s)}2. 
\end{equation}
We use \eqref{ae:exact} to define $E_{d,M}$ also for real $M$. This real
interpolation is necessary for discussing an exact collision parameter;
only the integer values are asserted to be graph Ehrhart polynomials.
The eventual collision parameters below are positive and much smaller
than $(d+1)/2$, so rounding them to integers gives legitimate graphs.

For clarity, the matching-deletion identity has a direct lattice-counting
proof. For an integer zero-sum vector $x$, put $q=\sum_i\max(x_i,0)$.
The complete graph gauge is $q$. For a complete multipartite graph the
gauge is
\[
 \max\left(q,\ \max_A\sum_{i\in A}|x_i|\right),
\]
where $A$ runs over its parts. Indeed, direct transport from positive to
negative coordinates has only the same-part restrictions. Hall's
deficiency is $\delta=\max_A(P_A+N_A-q)_+$: a supply set meeting two
parts reaches every demand, so only a single-part set can be deficient.
At most one part has positive deficiency. Send $q-\delta$ units directly
and send the remaining $\delta$ units within that part through an outside
vertex using two edges each. This gives cost $q+\delta$, and the displayed
lower bounds on cost are immediate. The argument also works with real
flows, so it computes the polytope gauge.

For a deleted matching the nontrivial parts are disjoint pairs. A point
of $kP_{K_{d+1}}$ is lost at a pair $\{1,2\}$ precisely when
$|x_1|+|x_2|>k$. It cannot be lost at two pairs, since
$\sum_i|x_i|=2q\le2k$. The two coordinates of a lost pair have opposite
signs. Fix the remaining vector $z\in\mathbb Z^{d-1}$ with positive and
negative masses $P,N$. For one choice of the two signs, write
$x_1=N+c$, $x_2=-P-c$ with $c\ge1$. The lost dilations satisfy
$P+N+c\le k\le P+N+2c-1$. Hence the generating function of the loss
from one pair is
\[
 2\sum_{z\in\mathbb Z^{d-1}}t^{\|z\|_1}
       \sum_{c\ge1}\frac{t^c-t^{2c}}{1-t}
 =\frac{2t(1+t)^{d-2}}{(1-t)^{d+1}}.
\]
This equals the Ehrhart series of $(C_d-C_{d-2})/2$. Additivity of the
disjoint losses proves \eqref{ae:exact} at every nonnegative integer
$s=k$, and therefore as a polynomial identity.

The preceding report established strict subcritical and supercritical
statements at the scale $M\asymp\sqrt d$, and left the equalities out.
The results here describe the previously unresolved local collision
windows. They are unconditional. They use no claim about RH or GRH.
They concern specified clusters of zeros near the canonical line
$\operatorname{Re}s=-1/2$. In particular, ``critical'' below means the
unique collision in that specified cluster, not a claim that this is the
first off-line zero anywhere in the polynomial.

Set
\begin{align}
 L&=\log d,&c_*&=\sqrt{8/\pi},&\beta&=2-\log 2,\label{ae:constants1}\\
 C&=\gamma+\frac53\log2+\frac23,&
 \kappa&=\frac{\pi^2}{4}-2+\frac23(2-\log2)^2.\label{ae:constants2}
\end{align}
Here $\gamma$ is Euler's constant. Numerically,
$C=2.3991276325\ldots$, $\kappa=1.6059772947\ldots$, and
$\kappa\pi^2=15.8503605761\ldots$.
Throughout, $h=(2j+1)\pi$ with a fixed integer $j\ge0$.
The fixed-$j$ condition is essential to the stated error estimates.

\begin{theorem}[The even-dimensional collision]\label{ae:thm-even}
Let $d$ tend to infinity through even integers, and suppose
\[
 M=c_*\sqrt d\left(1+\frac{\alpha_d}{L^2}\right),
 \qquad \alpha_d\longrightarrow\alpha\ne\kappa h^2.
\]
There are precisely two zeros in the corresponding local cluster, counted
with multiplicity, and they satisfy
\begin{equation}\label{ae:even-roots}
 s_\pm=-\frac12+\frac{ih}{L}
 +\frac{-ihC\ \pm\sqrt{\frac23(\alpha-\kappa h^2)}}{L^2}
 +o(L^{-2}).
\end{equation}
The square root may be taken in either order. If $\alpha<\kappa h^2$,
both roots are simple and lie exactly on the canonical line. If
$\alpha>\kappa h^2$, both are simple and off that line, forming a reflected
pair. Their complex conjugates form the corresponding negative-height
cluster.

More precisely, for every fixed sufficiently small neighborhood of
$\alpha=\kappa h^2$, there is for all large even $d$ a unique real local
collision parameter $M^{\mathrm{ev}}_{d,h}$ and collision height
$t^{\mathrm{ev}}_{d,h}$ with
\begin{align}
 M^{\mathrm{ev}}_{d,h}
 &=c_*\sqrt d\left(1+\frac{\kappa h^2}{L^2}+O(L^{-3})\right),
 \label{ae:even-critical}\\
 t^{\mathrm{ev}}_{d,h}
 &=\frac hL-\frac{hC}{L^2}+O(L^{-3}).\label{ae:even-height}
\end{align}
At that parameter the local root has multiplicity exactly two. In the
local neighborhood, increasing $M$ through this value changes two simple
on-line roots into a simple off-line reflected pair.
\end{theorem}

The meaning of ``precisely two in the local cluster'' in this and the next
theorem is unambiguous under rescaling: for any bounded parameter set
separated from its stated transition, choose a sufficiently large fixed
$R$. The cluster is the disk
$|L^2(s+1/2-ih/L)|<R$ in Theorem~\ref{ae:thm-even}, and the disk
$|L^{4/3}(s+1/2-ih/L)|<R$ in Theorem~\ref{ae:thm-odd}.
For all sufficiently large $d$ the indicated number of zeros lies in the
disk, all zeros there have the expansions given, and no zeros lie on its
boundary. A larger fixed $R$ gives the same cluster.

\begin{theorem}[The odd-dimensional cubic window]\label{ae:thm-odd}
Let $d$ tend to infinity through odd integers, and suppose
\[
 M=2c_*\sqrt d\left(1+\frac{\alpha_d}{L^{2/3}}\right),
 \qquad \alpha_d\longrightarrow\alpha.
\]
Define
\begin{equation}\label{ae:alpha-star}
 b=(\beta h)^{1/3},\qquad \alpha_*=-\frac32 b^2.
\end{equation}
If $\alpha\ne\alpha_*$, the cluster consists of exactly three simple
zeros, whose locations are
\begin{equation}\label{ae:odd-roots}
 s_k=-\frac12+\frac{ih}{L}+
             \frac{i x_k}{L^{4/3}}+o(L^{-4/3}),\qquad k=1,2,3,
\end{equation}
where $x_1,x_2,x_3$ are the roots of the explicit cubic
\begin{equation}\label{ae:cubic}
 x^3+2\alpha x-2\beta h=0.
\end{equation}
If $\alpha<\alpha_*$, all three exact zeros lie on the canonical line.
If $\alpha>\alpha_*$, exactly one does, and the other two form an off-line
reflected pair.

There is a unique local real collision parameter and height satisfying
\begin{align}
 M^{\mathrm{odd}}_{d,h}
 &=2c_*\sqrt d\left(1-
       \frac{3(\beta h)^{2/3}}{2L^{2/3}}+O(L^{-1})\right),
 \label{ae:odd-critical}\\
 t^{\mathrm{odd}}_{d,h}
 &=\frac hL-\frac{(\beta h)^{1/3}}{L^{4/3}}+O(L^{-5/3}).
 \label{ae:odd-height}
\end{align}
At this parameter two roots coalesce into a root of multiplicity exactly
two; the third remains simple, at scaled coordinate $x=2b+o(1)$.
In the local neighborhood, increasing $M$ through the collision value
changes three on-line roots into one on-line root and an off-line pair.
\end{theorem}

\begin{corollary}[What happens at the two previously omitted equalities]
\label{ae:cor-equality}
For even $d$, either nearest integer to $c_*\sqrt d$ gives two simple
on-line zeros in each of the fixed clusters in
Theorem~\ref{ae:thm-even}. For odd $d$, either nearest integer to
$2c_*\sqrt d$ gives an off-line pair in each of the fixed clusters in
Theorem~\ref{ae:thm-odd}. In the odd case put $r=(2\beta h)^{1/3}$.
The off-line pair has the explicit form
\[
 s_\pm=-\frac12\pm\frac{\sqrt3\,r}{2L^{4/3}}
       +i\left(\frac hL-\frac{r}{2L^{4/3}}\right)
       +o(L^{-4/3}),
\]
and the remaining zero is
$-1/2+i(h/L+r/L^{4/3})+o(L^{-4/3})$.
\end{corollary}

\subsection{Analytic input, including the resonance at the center}
The argument needs asymptotics in the dimension $d$, not a hypothesis on
arithmetic zeros. The elementary formulas
\begin{align}
 C_d(s)&=\sum_{k=0}^d 2^k\binom dk\binom sk,\label{ae:Cexact}\\
 H_d(s)&=\sum_{k=0}^d\binom dk\binom{d+k}k\binom sk\label{ae:Hexact}
\end{align}
lead to the dimension generating functions
\begin{align}
 \sum_{d\ge0}C_d(s)z^d
  &=(1+z)^s(1-z)^{-1-s},\label{ae:Cgen}\\
 \sum_{d\ge0}H_d(s)z^d
  &=(1-z)^{-1-2s}\,{}_2F_1(-s,-s;1;z^2).\label{ae:Hgen}
\end{align}
One can alternatively obtain \eqref{ae:Hexact} from the complete graph
numerator $\sum_k\binom dk^2 t^k$, followed by coefficient extraction in
its Ehrhart series. These are polynomial identities in $s$; checking the
lattice-counting generating functions for nonnegative integers $s$ proves
them for all complex $s$.

Write $s^*=-1-s$ and
\[
 A(s)=\Gamma(1+s)^{-2},\qquad B(s)=2^s\Gamma(s)^{-1},
 \qquad \epsilon=(-1)^d.
\]
Transfer at $z=1,-1$ in \eqref{ae:Cgen}--\eqref{ae:Hgen} gives
\begin{align}
 H_d(s)&=A(s)d^{2s}+\epsilon A(s^*)d^{2s^*}
             +O(d^{2\operatorname{Re}s-1}+d^{-3-2\operatorname{Re}s}),
 \label{ae:Hasy}\\
 D_d(s)&=B(s)d^{s-1}+\epsilon B(s^*)d^{s^*-1}
             +O(d^{\operatorname{Re}s-2}+d^{-3-\operatorname{Re}s}).
 \label{ae:Dasy}
\end{align}
These estimates are uniform on compact parameter sets avoiding the
resonances $1+2s\in\mathbb Z$; only the extension across $s=-1/2$ is
needed below. For completeness, near $z=1$ the hypergeometric connection
formula yields a singular term
\[
 (1-z)^{-1-2s}\frac{\Gamma(1+2s)}{\Gamma(1+s)^2}
                      (1+O(1-z))
\]
plus an analytic term. The corresponding singular term at $-1$ is
$(1+z)^{1+2s}\Gamma(-1-2s)/\Gamma(-s)^2$ times an analytic factor
with value one. Binomial coefficient transfer gives \eqref{ae:Hasy}.
Multiplying \eqref{ae:Cgen} by $(1-z^2)/2$ gives the generating function
for $D_d$, and the same transfer gives \eqref{ae:Dasy}. Local analytic
terms may be subtracted to any fixed order; both singular expansions and
their remainder bounds are uniform on a compact nonresonant parameter
set. This is the standard two-singularity Darboux argument, equivalently
multiple-singularity transfer~\cite[Chapter VI]{FS2009}; the hypergeometric
connection formulas are recorded in~\cite[Section 15.8]{DLMF}.

\begin{lemma}[Uniform holomorphic error at the resonant center]
\label{ae:lemma-uniform}
Fix $0<\eta<1/4$, a bounded set of $\mu\in\mathbb C$, and write
$s=-1/2+z$, $u=Lz$, $M=\mu\sqrt d$. Then throughout $|z|\le\eta$,
\begin{align}
 dE_{d,M}(-1/2+z)={}& A(-1/2+z)e^{2u}
       +\epsilon A(-1/2-z)e^{-2u}\notag\\
 &-\mu\{B(-1/2+z)e^u
       +\epsilon B(-1/2-z)e^{-u}\}
       +O_\eta(d^{-1+2\eta}).\label{ae:uniform}
\end{align}
On any fixed bounded $u$-disk, the same error bound holds after any fixed
number of $u$-derivatives. It also holds after one $\mu$-derivative;
higher such derivatives vanish.
\end{lemma}

\begin{proof}
The circle $|s+1/2|=\eta$ contains no resonance. On it,
\eqref{ae:Hasy}--\eqref{ae:Dasy}, multiplied by $d$, give the stated
absolute error uniformly for bounded $\mu$; the largest exponent is
$-1+2\eta$. The difference between the two sides of \eqref{ae:uniform}
without its error term is holomorphic in the whole disk. In particular,
the apparent resonance occurs only in the connection formula used to
estimate the boundary; it is not a singularity of that difference.
The maximum-modulus principle extends the bound to the disk. For all
large $d$, a fixed slightly larger $u$-disk maps strictly inside it.
Cauchy's derivative estimates then give all fixed-order $u$-derivative
bounds. The dependence on $\mu$ is affine, and its coefficient has the
same separately proved bound.
\end{proof}

\subsection{The two normal forms}
Define holomorphic functions near zero by
\[
 a(z)=\frac{A(-1/2+z)}{A(-1/2)},\qquad
 b(z)=\frac{B(-1/2+z)}{B(-1/2)}.
\]
Here $A(-1/2)=1/\pi$ and $c_*B(-1/2)=-1/\pi$. Write
\begin{align*}
 a(z)&=1+a_1z+a_2z^2+O(z^3),&
 a_1&=2\gamma+4\log2,&a_2&=\frac{a_1^2}{2}-\frac{\pi^2}{2},\\
 b(z)&=1+b_1z+b_2z^2+O(z^3),&
 b_1&=\gamma+3\log2-2,&b_2&=\frac{b_1^2}{2}-\frac{\pi^2}{4}-2.
\end{align*}
The values follow from the gamma recurrence and
$\psi(1/2)=-\gamma-2\log2$, $\psi'(1/2)=\pi^2/2$.
In particular
\begin{equation}\label{ae:algebra}
 a_1-2b_1=2\beta,\qquad
 \frac{2a_1-b_1}{3}=C,\qquad
 \frac{3C^2-2(a_2-b_2)}2=\kappa.
\end{equation}

For even $d$, put $u=ih+V/L$ and
$\mu=c_*(1+\alpha/L^2)$. Since $e^{ih}=-1$, divide
\eqref{ae:uniform} by $A(-1/2)$ and multiply by $L^2$.
Taylor expansion, using $z=u/L$, gives
\begin{equation}\label{ae:even-normal}
 \frac{L^2d}{A(-1/2)}E_{d,M}
 \left(-\frac12+\frac{ih}{L}+\frac V{L^2}\right)
 =3(V+ihC)^2-2(\alpha-\kappa h^2)+O(L^{-1}).
\end{equation}
The error is uniform on every fixed compact set of complex $(V,\alpha)$,
with all derivatives of any fixed order on smaller compact sets. Indeed,
the Taylor remainder after multiplication by $L^2$ is $O(L^{-1})$,
and the error from Lemma~\ref{ae:lemma-uniform} is
$O(L^2d^{-1+2\eta})=o(L^{-N})$ for every fixed $N$.
One direct way to check the coefficients is to expand first in
$v=u-ih$:
\[
 3v^2+(4a_1-2b_1)zv+2(a_2-b_2)z^2-2\alpha/L^2.
\]
Multiplication by $L^2$, followed by \eqref{ae:algebra}, gives
\eqref{ae:even-normal}.

For odd $d$, put $u=ih+i x/L^{1/3}$ and
$\mu=2c_*(1+\alpha/L^{2/3})$. The normalized leading expression is now
\[
 a(z)e^{2v}-a(-z)e^{-2v}
 -2(1+\alpha/L^{2/3})\{b(z)e^v-b(-z)e^{-v}\},\qquad v=u-ih.
\]
The linear terms in $v$ cancel. The remaining leading terms are
\[
 2v^3-4\alpha L^{-2/3}v+2(a_1-2b_1)z.
\]
Consequently
\begin{equation}\label{ae:odd-normal}
 -\frac{Ld}{2iA(-1/2)}E_{d,M}
 \left(-\frac12+\frac{ih}{L}+\frac {ix}{L^{4/3}}\right)
 =x^3+2\alpha x-2\beta h+O(L^{-1/3}).
\end{equation}
Again the error is holomorphic and locally uniform with every fixed
number of derivatives. Its largest term comes from retaining
$z=ih/L+ix/L^{4/3}$ in the last displayed expression. All other neglected
terms are of the indicated or smaller order. The dimension-asymptotic
error is exponentially small compared with these logarithmic scales.

\subsection{Proof of the collision theorems}
\begin{proof}[Proof of Theorem~\ref{ae:thm-even}]
If $\alpha\ne\kappa h^2$, the limiting quadratic in
\eqref{ae:even-normal} has two distinct roots. Rouch\'e's theorem on
small disjoint circles around them, and on a sufficiently large fixed
circle enclosing both, proves the root count, simplicity, and
\eqref{ae:even-roots}. Reflection across the canonical line preserves
the exact zero multiset, by real coefficients and
$E_{d,M}(-1-s)=(-1)^dE_{d,M}(s)$. When the limiting roots are on the
line, choose reflection-invariant disks centered on that line. Each
contains exactly one exact root counted with multiplicity, so the root
is exactly on-line and simple. When the limiting roots have nonzero
real offsets, their exact roots are off-line for all large $d$.

For the exact collision, set $V=ix$ in \eqref{ae:even-normal}. The
result is real for real $x,\alpha$, and tends in every fixed derivative
norm to
\[
 q(x,\alpha)=-3(x+hC)^2+2(\kappa h^2-\alpha).
\]
The equations $q=q_x=0$ have the solution
$(x,\alpha)=(-hC,\kappa h^2)$. The Jacobian of $(q,q_x)$ with
respect to $(x,\alpha)$ there is
$\left(\begin{smallmatrix}0&-2\\-6&0\end{smallmatrix}\right)$,
which is invertible. The quantitative implicit-function theorem,
or its proof by a contraction on a fixed small neighborhood, gives a
unique nearby solution for the exact functions, with errors $O(L^{-1})$.
This proves \eqref{ae:even-critical}--\eqref{ae:even-height}.
The second $x$-derivative and the $\alpha$-derivative stay strictly
negative. The local graph therefore has a strict nondegenerate maximum,
which crosses zero as $M$ increases. It has two real roots on the lower
side, a double real root at the crossing, and no real roots on the upper
side. The two-root complex count and reflection then give the claimed
off-line pair. The nonzero second derivative also proves multiplicity
exactly two.
\end{proof}

\begin{proof}[Proof of Theorem~\ref{ae:thm-odd}]
The discriminant of \eqref{ae:cubic} is
\[
 -32\alpha^3-108\beta^2h^2.
\]
It is positive, respectively negative, precisely for
$\alpha<\alpha_*$, respectively $\alpha>\alpha_*$. Thus the cubic
has three distinct real roots on the first side, and one real root and
a nonreal conjugate pair on the second. Apply Rouch\'e to
\eqref{ae:odd-normal} as in the even case. In the $x$ coordinate,
reflection across the canonical line is complex conjugation, and the
normalized exact function is real for real $x,\alpha$. The simple-root
uniqueness argument therefore proves exact on-line location, not merely
an asymptotically vanishing displacement, for each limiting real root.

At $\alpha=\alpha_*$, the cubic is $(x+b)^2(x-2b)$. At the double
root $x=-b$, the Jacobian of $(q,q_x)$ is
\[
 \begin{pmatrix}0&-2b\\-6b&2\end{pmatrix},\qquad
 \det=-12b^2\ne0.
\]
The implicit-function theorem gives a unique nearby exact collision
with $(x,\alpha)=(-b,\alpha_*)+O(L^{-1/3})$. This yields
\eqref{ae:odd-critical}--\eqref{ae:odd-height}. Both $q_\alpha$ and
$q_{xx}$ are negative at the double root, so the same local-maximum
argument proves the direction of the transition. The third root stays
simple and real near $2b$. The nonzero second derivative at the colliding
pair gives multiplicity exactly two in the exact finite-$d$ polynomial.
\end{proof}

\begin{proof}[Proof of Corollary~\ref{ae:cor-equality}]
Rounding changes $M$ by $O(1)$. Hence the even rescaled parameter changes
by $O(L^2/\sqrt d)=o(1)$ and the odd one by
$O(L^{2/3}/\sqrt d)=o(1)$. Both limits are therefore $\alpha=0$.
Since $\kappa>0$, this is on the even on-line side; since
$\alpha_*<0$, it is on the odd off-line side. At $\alpha=0$ the cubic
has roots $r,r e^{2\pi i/3},r e^{-2\pi i/3}$, giving the stated formulas.
\end{proof}

\subsection{Exact checks and the limits of the result}
The supplementary programs check the exact polynomial identity, reflection,
and both dimension recurrences by independent rational arithmetic for
small dimensions. They check the Taylor coefficient identities symbolically
and compute local collision parameters with high-precision arithmetic
using the stable real recurrences
\[
 (d+1)^2h_{d+1}=2(2d+1)t h_d-d^2h_{d-1},\qquad
 (d+1)c_{d+1}=2t c_d-dc_{d-1},
\]
where $h_d(t)=i^{-d}H_d(-1/2+it)$ and
$c_d(t)=i^{-d}C_d(-1/2+it)$. These computations illustrate the
convergence; the asymptotic theorems are proved above independently of
the computations.

The result determines local multiplicities, displacement scales, and
collision thresholds for every fixed odd multiple $h$ of $\pi$.
It does not establish the complete bipartite all-width canonical-line
conjecture, a global finite-dimensional graph threshold, a statement
uniform in growing $j$, or an exact critical parameter for integer
$M$ when a real collision lies between consecutive integers. It is a
human-readable mathematical proof, not a Lean formalization. Its novelty
claim is relative to the inspected repository and incoming report;
publication priority has not been exhaustively established.
